\section{The simulator: implementation, information safety, verification}
\label{sec:engine}

\subsection{Implementation}

The engine is a single-header C++20 core in \filepath{engine/src/}. A hand is one
\verb|uint64_t| over the 54 card indices $c = 6h + i$, where $h$ numbers the nine
half-suits and $i$ the six cards within one, so a half-suit mask is
\verb|0x3F << 6h| and the six hands together with the deal they came from occupy
under a hundred bytes. Where the policy explores a hypothetical --- the two-ply
refinement of \S\ref{sec:twoply} --- it copies the constraint state and
recomputes on the copy rather than making and unmaking moves on the live game.
There are no rollouts anywhere in this agent: every quantity the policy consumes
is either a closed-form posterior query (\S\ref{sec:belief}) or a fitted linear
score (\S\ref{sec:policy}).

Games are driven by a loop that, after every public event, (i) offers every
player the opportunity to declare, (ii) resolves a cardless team into the forced
endgame, (iii) resolves a cardless turn-holder into a chosen successor, and
(iv) otherwise asks the turn-holder for a move. All randomness comes from
counter-based streams keyed by \emph{(seed, seat, purpose)}, so a deal and its
entire play are reproducible from the seed alone. The duplicate-block protocol
of \S\ref{sec:protocol} depends on that property.

\subsection{Information safety}

The acting policy never receives the game state. It receives its own
\texttt{Knowledge} object --- the constraint bookkeeping of
\S\ref{sec:belief} --- and the public record, in which each event carries only
the actor, the target, the named card, the outcome, any declared allocation, and
the resulting public hand counts. There is no interface through which a policy
can read another player's cards. A leak of the kind the \vthree{} audit reported
in its predecessor \cite{fishbot03}, where a reply-risk feature estimated a
target's follow-up options by invoking the legal-move generator with that target
as actor, is therefore excluded structurally rather than by inspection.

To confirm that no leak remains, and more importantly that the deductions the
belief machinery makes are \emph{sound}, the engine has an audit mode that runs
after every public event and checks, for every player, four properties against
the true deal:

\begin{itemize}
  \item every card the constraint state records as resolved is in fact held by
        the seat it is resolved to, and no card of a live half-suit is marked out
        of play;
  \item for every unresolved card, the true holder lies in the surviving support
        $M_c$ (\S\ref{sec:belief});
  \item the derived capacities \eqref{eq:capacity} equal the true counts of
        unresolved cards per player;
  \item every live ask-legality certificate \eqref{eq:disj} is satisfied by the
        true deal.
\end{itemize}

A single violation would mean the belief engine had excluded the truth. In
experiment E1 the engine performed \textbf{\numAuditChecks} such checks with
\textbf{\numAuditViolations} violations under the full dialect, and
\textbf{\numAuditChecksLegacy} checks with no violations under the \vthree{}
dialect (\S\ref{sec:results-verify}). The audit is run inside \texttt{fish
verify}, which plays the $9 \times 9$ ordered cross-product of nine policies
(\vfast{}, \vthree{}, \vtwo{}, turn-starvation lockout, posterior detective,
focused hunter, adaptive diversifier, misdirection artist, random legal control),
so that both orderings of every pair and every self-play pairing are exercised.

\subsection{Verification suite}

Beyond the soundness audit the suite checks deck integrity and card
conservation; that exactly nine half-suits are resolved in every completed game;
that a declaration only ever assigns cards to the declarer's own teammates; that
ask legality is enforced independently on the move generator and on the applier,
so a policy cannot submit an illegal ask even if its generator is wrong; that
replaying a seed at a fixed thread count reproduces the game bit-identically;
and that the action cap of \S\ref{sec:rules} --- a backstop whose residue
adjudication is a convention of the simulator rather than a rule of the game ---
is not reached in ordinary play, which it was not (\numLimitRate{}, E1 and E3).
The in-policy mechanism that keeps games terminating is the subject of
\S\ref{sec:termination}. The random streams are keyed so that thread count
cannot affect a game's outcome, but the suite tests replay identity only at a
fixed thread count and does not test that claim directly.
